\documentclass{article}
\usepackage{graphicx} % Required for inserting images
\usepackage{amsmath} %USAR ESSA BIBLIOTECAA!!!!

\section{Definição da Amostra}

Para a definição do conjunto de observações a ser utilizado nas análises
estatísticas subsequentes, considera-se o maior número de matrícula entre
os integrantes do grupo. Os números de matrícula são apresentados na
Tabela~\ref{tab:matriculas}.

\begin{table}[h]
    \centering
    \begin{tabular}{lc}
        \hline
        \textbf{Integrante} & \textbf{Matrícula} \\
        \hline
        Allan Martins Gadelha & 579087 \\
        Clarisse Maria Cabral Montenegro & 580585 \\
        Maria Luisa Fernandes de Mendonça & 581706 \\
        Sara de Sousa Magalhães & 581189 \\
        \hline
    \end{tabular}
    \caption{Números de matrícula dos integrantes do grupo.}
    \label{tab:matriculas}
\end{table}

A maior matrícula entre os integrantes é:

\[
M = 581706.
\]

Conforme a definição estabelecida, calcula-se o valor de
$r$ por meio da expressão:

\[
r = 1 + (M\mod 100).
\]

O operador $\mod$ representa o resto da divisão inteira de $M$ por 100.

Calculando o resto da divisão de $M$ por 100, obtém-se:

\[
581706 \text{ mod } 100 = 6.
\]

Logo,

\[
r = 1 + 6 = \boxed{7}.
\]

Portanto, a seleção das observações deve ser iniciada na observação de
número 7. Considerando a necessidade de selecionar 300 observações
consecutivas, a última observação selecionada é determinada por:

\[
r + 300 - 1 = 7 + 300 - 1 = 306.
\]

Dessa forma, a amostra utilizada nas análises estatísticas subsequentes
é composta pelas 300 observações compreendidas entre as observações 7 e
306, inclusive:

\[
\boxed{\{7,8,9,\ldots,306\}}.
\]

